\documentclass[11pt,a4paper]{article}
\usepackage[T1]{fontenc}
\usepackage[utf8]{inputenc}
\usepackage[french]{babel}
\usepackage{lmodern}
\usepackage{amsmath,amssymb,amsthm,mathtools}
\usepackage[margin=2.35cm]{geometry}
\usepackage{booktabs,tabularx}
\usepackage{enumitem}
\usepackage{xcolor}
\usepackage{microtype}
\usepackage{hyperref}

\hypersetup{colorlinks=true,linkcolor=blue!45!black,urlcolor=blue!45!black}
\setlist{nosep}
\newcommand{\E}{\mathbb E}
\newcommand{\Pp}{\mathbb P}
\newcommand{\R}{\mathbb R}
\newcommand{\1}{\mathbf 1}
\newcommand{\dd}{\,\mathrm d}
\newcommand{\Law}{\mathcal L}
\newcommand{\cim}{\dagger}
\DeclareMathOperator{\supp}{supp}

\newtheorem{theorem}{Théorème}[section]
\newtheorem{proposition}[theorem]{Proposition}
\newtheorem{lemma}[theorem]{Lemme}
\newtheorem{corollary}[theorem]{Corollaire}
\theoremstyle{definition}
\newtheorem{definition}[theorem]{Définition}
\theoremstyle{remark}
\newtheorem{remark}[theorem]{Remarque}

\title{Étude analytique d'un prêt unique entre deux entités\\
\large Effets du principal, du taux, de la concavité, de la dépréciation
et des chocs lognormaux}
\author{Complément mathématique à la spécification M4.2}
\date{29 juillet 2026}

\begin{document}
\maketitle

\begin{abstract}
On isole deux entités initialement sans contrat. Au pas initial, la plus
riche accorde à la plus pauvre un unique prêt nominal perpétuel ; aucun
autre prêt et aucune naissance n'ont ensuite lieu. La production est
$K\mapsto K^\gamma$, la dépréciation vaut $\delta$, et les chocs sont
exactement ceux de M4.2. Pour un principal général $q$ et un taux
$r\in[0,1]$, nous dérivons la dynamique tuée de l'emprunteuse, les noyaux
intégraux donnant la loi exacte de son temps de défaut et la loi jointe des
deux capitaux au défaut. Deux règles sont comparées : égalisation à la
moyenne arithmétique et cible à la moyenne géométrique ; la première
maximise la production initiale mais donne un temps de survie presque
sûrement inférieur ou égal. Deux conclusions structurelles apparaissent.
Premièrement, la prêteuse ne peut jamais
mourir dans ce système orienté à deux sommets. Deuxièmement, si le choc est
non dégénéré, l'emprunteuse meurt presque sûrement en temps fini, tandis
que le capital de la prêteuse converge en loi vers la distribution
invariante de la dynamique autarcique. Le régime déterministe est traité
séparément : lui seul admet des trajectoires convergeant vers deux
constantes positives.
\end{abstract}

\tableofcontents

\section{Cadre exact et convention temporelle}

On fixe
\[
0<\gamma<1,\qquad 0\leq\delta<1,\qquad d=1-\delta>0,
\qquad \sigma\geq0,
\]
et, conformément au modèle principal M4.2, $A=1$. Posons
\begin{equation}
\Phi_\gamma(z)=z+z^\gamma,\qquad z\geq0.
\label{eq:phi}
\end{equation}
Cette fonction regroupe le capital après choc et la production retenue,
avant service des intérêts et dépréciation. Elle est continue, strictement
croissante et bijective de $\R_+$ sur $\R_+$. Son inverse est noté
$\Psi_\gamma=\Phi_\gamma^{-1}$.

Les deux capitaux juste avant l'unique transfert sont
\[
k_1>k_2>0.
\]
L'entité 1 prête un principal $q\in(0,k_1)$ à l'entité 2, au taux fixe
$r\in[0,1]$. Le service nominal perpétuel est
\begin{equation}
c=rq.
\label{eq:c}
\end{equation}
Le temps $n=0$ est placé immédiatement après ce transfert :
\begin{equation}
L_0=k_1-q,\qquad B_0=k_2+q,
\label{eq:initial}
\end{equation}
où $L$ désigne le capital de la prêteuse et $B$ celui de l'emprunteuse.
Cette convention sépare l'effet du contrat des phases choc--production du
pas où il est émis. Le transfert conserve le capital réel total et les
valeurs nettes initiales :
\[
(L_0+q)=k_1,\qquad (B_0-q)=k_2.
\]

À chaque pas $n+1$, indépendamment entre entités et entre dates,
\begin{equation}
\xi^i_{n+1}\sim\mathcal N(-\sigma^2/2,\sigma^2),\qquad
X^i_{n+1}=e^{\xi^i_{n+1}},\qquad \E[X^i_{n+1}]=1.
\label{eq:shock}
\end{equation}
Pour $\sigma>0$, la densité du multiplicateur est
\begin{equation}
\rho_\sigma(x)=\frac{1}{x\sigma\sqrt{2\pi}}
\exp\!\left[-\frac{(\log x+\sigma^2/2)^2}{2\sigma^2}\right],
\qquad x>0.
\label{eq:lognormal-density}
\end{equation}
Le cas $\sigma=0$ correspond à $X\equiv1$.

\paragraph{Portée des tolérances.}
Les résultats sont écrits pour les seuils mathématiques exacts. La
tolérance d'insolvabilité $10^{-9}$ de l'implémentation remplace simplement
$q$ par $q-10^{-9}$ dans le test de valeur nette ; elle ne modifie aucune
conclusion structurelle tant que les capitaux et le prêt sont à une
échelle macroscopique.

\section{Deux définitions du principal}

\subsection{Égalisation à la moyenne arithmétique}

Imposer $L_0=B_0=(k_1+k_2)/2$ donne
\begin{equation}
q_{\mathrm A}=\frac{k_1-k_2}{2},
\qquad L_0^{\mathrm A}=B_0^{\mathrm A}=\frac{k_1+k_2}{2}.
\label{eq:qA}
\end{equation}
Cette règle est l'unique maximisatrice de la production immédiatement
postérieure au transfert. En effet, sur $0\leq q\leq k_1-k_2$,
\[
q\longmapsto (k_1-q)^\gamma+(k_2+q)^\gamma
\]
est strictement concave et sa dérivée s'annule uniquement lorsque les deux
capitaux sont égaux. C'est la version à deux dimensions du principe
d'égalisation d'une fonction symétrique strictement concave sous contrainte
de somme.

\subsection{Cible à la moyenne géométrique}

La règle de M4.2 vise pour l'emprunteuse le capital
$g=\sqrt{k_1k_2}$. Le principal général de la spécification est
\[
\min(k_1-g,g-k_2).
\]
Or
\[
(k_1-g)-(g-k_2)=k_1+k_2-2\sqrt{k_1k_2}
=(\sqrt{k_1}-\sqrt{k_2})^2>0.
\]
Par conséquent,
\begin{equation}
q_{\mathrm G}=\sqrt{k_1k_2}-k_2,\qquad
B_0^{\mathrm G}=\sqrt{k_1k_2},\qquad
L_0^{\mathrm G}=k_1+k_2-\sqrt{k_1k_2}.
\label{eq:qG}
\end{equation}

\begin{proposition}[Ordre des deux transferts]
Pour $k_1>k_2>0$,
\[
0<q_{\mathrm G}<q_{\mathrm A},\qquad
B_0^{\mathrm G}<B_0^{\mathrm A},\qquad
L_0^{\mathrm G}>L_0^{\mathrm A}.
\]
De plus, le levier initial $q/B_0$ est plus grand sous la règle
arithmétique :
\begin{equation}
\frac{q_{\mathrm A}}{B_0^{\mathrm A}}
>\frac{q_{\mathrm G}}{B_0^{\mathrm G}}.
\label{eq:leverage-order}
\end{equation}
\end{proposition}
\begin{proof}
La première différence vaut
\[
q_{\mathrm A}-q_{\mathrm G}
=\frac{(\sqrt{k_1}-\sqrt{k_2})^2}{2}>0.
\]
Pour le levier, avec $s=\sqrt{k_1/k_2}>1$,
\[
\frac{q_{\mathrm A}}{B_0^{\mathrm A}}
-\frac{q_{\mathrm G}}{B_0^{\mathrm G}}
=\frac{s^2-1}{s^2+1}-\frac{s-1}{s}
=\frac{(s-1)^2}{s(s^2+1)}>0.
\]
\end{proof}

La règle arithmétique maximise donc la production commune instantanée,
mais crée simultanément un principal nominal et un levier plus élevés.
Après normalisation par le principal, ce second effet domine au sens
trajectoriel : on démontrera en proposition~\ref{prop:A-before-G} que
l'emprunteuse arithmétique meurt toujours au plus tôt aussi vite que
l'emprunteuse géométrique sous les mêmes chocs.

\section{Dynamique exacte tant que l'emprunteuse vit}

Supposons l'emprunteuse vivante au temps $n$. Après choc et production,
les ressources des deux entités sont
\[
Y^L_{n+1}=\Phi_\gamma(L_nX^L_{n+1}),\qquad
Y^B_{n+1}=\Phi_\gamma(B_nX^B_{n+1}).
\]
Le paiement effectif est
\begin{equation}
p_{n+1}=\min(c,Y^B_{n+1}).
\label{eq:payment}
\end{equation}
Juste avant la résolution des faillites, les capitaux sont
\begin{equation}
\widetilde L_{n+1}=d(Y^L_{n+1}+p_{n+1}),\qquad
\widetilde B_{n+1}=d(Y^B_{n+1}-p_{n+1}).
\label{eq:tentative}
\end{equation}

L'emprunteuse survit si elle paie tout l'intérêt et si sa valeur nette
post-dépréciation est non négative :
\[
Y^B_{n+1}\geq c,\qquad
d(Y^B_{n+1}-c)-q\geq0.
\]
Comme $q>0$, la seconde condition implique la première. On obtient donc le
critère unique
\begin{equation}
\boxed{\quad
Y^B_{n+1}\geq h_{q,r,\delta},\qquad
h_{q,r,\delta}=rq+\frac{q}{1-\delta}
=q\left(r+\frac1d\right).
\quad}
\label{eq:survival-threshold}
\end{equation}
Définissons le seuil en capital choqué
\begin{equation}
z_*=\Psi_\gamma(h_{q,r,\delta}).
\label{eq:zstar}
\end{equation}
Conditionnellement à $B_n=b$, le décès au pas suivant équivaut à
\begin{equation}
X^B_{n+1}<\frac{z_*}{b}
\quad\Longleftrightarrow\quad
\xi^B_{n+1}<\log\frac{z_*}{b}.
\label{eq:death-event}
\end{equation}

Sur l'événement de survie, la dynamique est ainsi
\begin{equation}
\begin{split}
L_{n+1}&=d\{\Phi_\gamma(L_nX^L_{n+1})+c\},\\
B_{n+1}&=d\{\Phi_\gamma(B_nX^B_{n+1})-c\}\geq q.
\end{split}
\label{eq:alive-dynamics}
\end{equation}
Au décès de l'emprunteuse, son capital résiduel pré-liquidation est
\begin{equation}
B_{T}^{\cim}=d\,[Y^B_T-c]_+\in[0,q),
\label{eq:terminal-borrower}
\end{equation}
puis ce capital est détruit et la créance $q$ est annulée. La taille
post-liquidation est donc le point cimetière, noté $\partial$ ; conserver
$B_T^{\cim}$ permet de mesurer ce qui est effectivement détruit.

\subsection{Risque conditionnel à un pas}

Notons $\Phi_{\mathrm N}$ la fonction de répartition normale standard. De
\eqref{eq:death-event}, pour $\sigma>0$,
\begin{equation}
\boxed{\quad
\kappa_{q,r}(b)
=\Pp(T=n+1\mid T>n,B_n=b)
=\Phi_{\mathrm N}\!\left(
\frac{\log(z_*/b)+\sigma^2/2}{\sigma}
\right).
\quad}
\label{eq:hazard}
\end{equation}
Le défaut de liquidité strict correspond au seuil plus bas
$z_{\mathrm{liq}}=\Psi_\gamma(c)$ :
\begin{equation}
\Pp(\text{liquidité au pas suivant}\mid B_n=b)
=\Phi_{\mathrm N}\!\left(
\frac{\log(z_{\mathrm{liq}}/b)+\sigma^2/2}{\sigma}
\right),
\label{eq:liquidity-hazard}
\end{equation}
avec la convention que cette probabilité vaut zéro si $c=0$. La
différence entre \eqref{eq:hazard} et \eqref{eq:liquidity-hazard} est la
probabilité d'une insolvabilité malgré paiement intégral de la rente.

Pour chacune des règles $s\in\{\mathrm A,\mathrm G\}$, les expressions
exactes au premier pas s'obtiennent en substituant
\begin{equation}
q=q_s,\qquad b=B_0^s,\qquad
z_*^s=\Psi_\gamma\!\left[q_s\left(r+\frac1d\right)\right]
\label{eq:schemes-substitution}
\end{equation}
dans \eqref{eq:hazard}. La règle arithmétique augmente à la fois $q$ et
$b$, mais l'augmentation relative du principal est la plus forte. Le
classement exact est
\begin{equation}
\kappa^{\mathrm A}_{q,r}(B_0^{\mathrm A})
>\kappa^{\mathrm G}_{q,r}(B_0^{\mathrm G}),
\label{eq:hazard-scheme-order}
\end{equation}
comme conséquence du résultat trajectoriel suivant.

\begin{proposition}[La règle arithmétique réduit la durée de vie]
\label{prop:A-before-G}
Couplons les deux règles avec la même suite de chocs de l'emprunteuse et
le même taux $r$. Alors
\begin{equation}
T_B^{\mathrm A}\leq T_B^{\mathrm G}
\qquad\text{presque sûrement}.
\label{eq:TA-TG}
\end{equation}
Pour $\sigma>0$, l'inégalité des risques au premier pas est stricte.
\end{proposition}
\begin{proof}
Normalisons le capital de l'emprunteuse par son principal : $W_n=B_n/q$.
Tant qu'elle vit, \eqref{eq:alive-dynamics} devient
\begin{equation}
W_{n+1}=d\left\{W_nX_{n+1}
+q^{\gamma-1}(W_nX_{n+1})^\gamma-r\right\},
\qquad W_{n+1}\geq1.
\label{eq:normalized-borrower}
\end{equation}
La proposition sur les leviers donne
$W_0^{\mathrm A}<W_0^{\mathrm G}$. De plus,
$q_{\mathrm A}>q_{\mathrm G}$ et $\gamma-1<0$, donc
$q_{\mathrm A}^{\gamma-1}<q_{\mathrm G}^{\gamma-1}$. Le membre de droite
de \eqref{eq:normalized-borrower} est croissant en $W$. Une récurrence
immédiate donne ainsi $W_n^{\mathrm A}\leq W_n^{\mathrm G}$ tant que la
première chaîne n'est pas absorbée. Le franchissement de $1$ se produit
donc d'abord sous la règle arithmétique. La densité lognormale est
strictement positive et les deux seuils sont distincts, d'où la stricte
inégalité au premier pas.
\end{proof}

\section{Temps d'arrêt et impossibilité d'une seconde mort}

On définit, dans $\mathbb N^*=\{1,2,\ldots\}$ complété par $+\infty$,
\begin{equation}
T_B=\inf\{n\geq1:B_n\text{ fait défaut au pas }n\},
\qquad
T_L=\inf\{n\geq1:L_n\text{ meurt au pas }n\}.
\label{eq:stopping-times}
\end{equation}

\begin{theorem}[La prêteuse est immortelle]
Pour tout $k_1>k_2>0$, tout principal admissible $0<q<k_1$, tout
$r\in[0,1]$, $0\leq\delta<1$ et toute réalisation finie des chocs,
\begin{equation}
T_L=+\infty.
\label{eq:lender-immortal}
\end{equation}
Ainsi, dans le modèle exact à deux entités,
\[
T_{(1)}=\min(T_B,T_L)=T_B,
\qquad T_{(2)}=\max(T_B,T_L)=+\infty.
\]
Une séquence de morts $T_1<T_2<\infty$ est impossible.
\end{theorem}
\begin{proof}
Avant la mort de l'emprunteuse, la valeur nette de la prêteuse est
$L_n+q>0$ : elle a une créance $q$, aucune dette, et son capital reste
strictement positif car $d>0$, $X^L_n>0$ et
$\Phi_\gamma(z)>0$ pour $z>0$. Au pas où l'emprunteuse meurt,
$\widetilde L=d\{\Phi_\gamma(LX^L)+p\}>0$. La perte de la créance retire
$q$ de l'actif nominal mais ne retire aucun capital réel ; la nouvelle
valeur nette est donc $\widetilde L>0$. Après l'annulation du contrat, la
prêteuse n'a ni créance ni dette et suit une dynamique autarcique qui
préserve strictement la positivité. Aucun critère de liquidité ou
d'insolvabilité ne peut être déclenché.
\end{proof}

Cette proposition n'est pas une approximation asymptotique : elle découle
de l'orientation du bilan. Une cascade à deux morts exigerait que la
prêteuse ait par ailleurs une dette ou une obligation de service, ce qui
est exclu par l'expérience présente.

\section{Loi exacte du temps de mort de l'emprunteuse}

Les formules suivantes caractérisent entièrement les distributions
demandées. Elles ne sont généralement pas réductibles à des fonctions
élémentaires, car l'itération non linéaire $z+z^\gamma$ est composée avec
un choc lognormal. Elles sont néanmoins exactes : elles ramènent le
problème à des intégrales unidimensionnelles positives.

\subsection{Noyau sous-stochastique de survie}

Pour $b\geq q$ et $y\geq q$, posons
\begin{equation}
u_y=\Psi_\gamma\!\left(c+\frac yd\right),
\qquad x(b,y)=\frac{u_y}{b}.
\label{eq:change-borrower}
\end{equation}
Le changement de variable
$y=d\{\Phi_\gamma(bx)-c\}$ donne le noyau de survie
\begin{equation}
Q_{q,r}(b,\dd y)=q_{q,r}(b,y)\1_{\{y\geq q\}}\dd y,
\label{eq:Q}
\end{equation}
avec
\begin{equation}
q_{q,r}(b,y)=
\frac{\rho_\sigma(x(b,y))}
{d\,b\,\Phi_\gamma'(u_y)},
\qquad
\Phi_\gamma'(u)=1+\gamma u^{\gamma-1}.
\label{eq:Qdensity}
\end{equation}
Sa masse est la probabilité de survie conditionnelle :
\begin{equation}
Q_{q,r}\1(b)=1-\kappa_{q,r}(b).
\label{eq:Qmass}
\end{equation}

Soit $\mu_0=\delta_{B_0}$ la masse de Dirac au capital initial. La mesure
non normalisée des survivantes au temps $n$ est
\begin{equation}
\mu_n=\mu_0Q_{q,r}^{,n}.
\label{eq:mun}
\end{equation}
Elle satisfait, pour $n\geq0$,
\begin{align}
\Pp(T_B>n)&=\mu_n\1,
\label{eq:survival-law}\\
\Pp(T_B=n+1)&=\int_q^\infty \kappa_{q,r}(b)\,\mu_n(\dd b).
\label{eq:Tlaw}
\end{align}
Les équations \eqref{eq:Qdensity}--\eqref{eq:Tlaw} donnent la loi discrète
exacte de $T_B$ en fonction de
$(k_1,k_2,\gamma,\delta,\sigma,r)$ et de la règle choisie pour $q$.

Sous forme récursive, si $m_n$ est la densité de $\mu_n$ pour $n\geq1$,
\begin{equation}
m_{n+1}(y)=\1_{\{y\geq q\}}
\int_q^\infty m_n(b)q_{q,r}(b,y)\dd b,
\label{eq:density-recursion}
\end{equation}
la première étape étant
$m_1(y)=q_{q,r}(B_0,y)\1_{\{y\geq q\}}$.

\subsection{Loi du capital détruit au temps d'arrêt}

Conditionnellement à $B_n=b$, la loi de $B_{n+1}^{\cim}$ sur l'événement
de mort comporte une masse en zéro et une densité sur $(0,q)$. La masse en
zéro, correspondant au défaut de liquidité, est
\begin{equation}
a_0(b)=
\Phi_{\mathrm N}\!\left(
\frac{\log(\Psi_\gamma(c)/b)+\sigma^2/2}{\sigma}
\right)
\label{eq:atom-zero}
\end{equation}
pour $c>0$, et $a_0(b)=0$ pour $c=0$. Pour $0<y<q$, la densité est
\begin{equation}
d_{q,r}(b,y)=
\frac{\rho_\sigma\!\left(
\Psi_\gamma(c+y/d)/b\right)}
{d\,b\,\Phi_\gamma'\!\left(\Psi_\gamma(c+y/d)\right)}.
\label{eq:death-density}
\end{equation}
Notons $D_{q,r}(b,\dd y)$ la mesure ainsi formée. Sa masse totale vérifie
\[
D_{q,r}\1(b)=a_0(b)+\int_0^q d_{q,r}(b,y)\dd y
=\kappa_{q,r}(b).
\]
La loi jointe du temps de mort et du capital pré-liquidation est alors
\begin{equation}
\boxed{\quad
\Pp(T_B=n+1,B_T^{\cim}\in\dd y)
=\int_q^\infty D_{q,r}(b,\dd y)\,\mu_n(\dd b).
\quad}
\label{eq:joint-T-B}
\end{equation}
En sommant sur $n$, la loi du capital détruit est la mesure-résolvante
\begin{equation}
\Law(B_T^{\cim})
=\delta_{B_0}(I-Q_{q,r})^{-1}D_{q,r}
=\sum_{n\geq0}\delta_{B_0}Q_{q,r}^{,n}D_{q,r},
\label{eq:resolvent-B}
\end{equation}
dès que l'absorption est presque sûre, propriété démontrée en
section~\ref{sec:asymptotic-stochastic}.

\section{Loi jointe des deux capitaux au premier décès}

La loi de $L_{T_B}$ ne se déduit pas du seul noyau $Q$, car le paiement
reçu au pas de mort dépend du choc de l'emprunteuse. Cette dépendance
disparaît lors d'une insolvabilité avec service complet, mais subsiste lors
d'un défaut de liquidité.

Pour toute fonction mesurable positive $g$, définissons le noyau de mort
joint
\begin{multline}
(\mathcal Dg)(l,b)=
\int_0^{z_*/b}\int_0^\infty
g\Bigl(
d\{\Phi_\gamma(lx_1)+\min(c,\Phi_\gamma(bx_2))\},\\
d[\Phi_\gamma(bx_2)-c]_+
\Bigr)
\rho_\sigma(x_1)\rho_\sigma(x_2)\dd x_1\dd x_2.
\label{eq:joint-death-kernel}
\end{multline}
Sur un pas où les deux entités continuent, les transitions sont
indépendantes conditionnellement à l'état courant, puisque le paiement est
alors exactement $c$. Le noyau de la prêteuse avec rente est, pour
$l'>dc$,
\begin{equation}
R_c(l,l')=
\frac{\rho_\sigma\!\left(
\Psi_\gamma(l'/d-c)/l\right)}
{d\,l\,\Phi_\gamma'\!\left(\Psi_\gamma(l'/d-c)\right)}.
\label{eq:lender-rent-kernel}
\end{equation}
Le noyau joint sous-stochastique de survie vaut donc
\begin{equation}
\mathcal S((l,b),\dd l'\dd b')
=R_c(l,l')q_{q,r}(b,b')
\1_{\{l'>dc,b'\geq q\}}\dd l'\dd b'.
\label{eq:joint-survival-kernel}
\end{equation}
Avec $\nu_0=\delta_{(L_0,B_0)}$, on obtient enfin
\begin{equation}
\boxed{\quad
\Pp(T_B=n+1,(L_{T_B},B_{T_B}^{\cim})\in A)
=(\nu_0\mathcal S^n\mathcal D)(A).
\quad}
\label{eq:joint-stopped-law}
\end{equation}
La somme
\begin{equation}
\Law(L_{T_B},B_{T_B}^{\cim})
=\nu_0(I-\mathcal S)^{-1}\mathcal D
\label{eq:joint-resolvent}
\end{equation}
est la réponse analytique complète à la question des tailles au premier
temps d'arrêt. Au second temps d'arrêt, $T_L=+\infty$ ; il n'existe donc
pas de variable $L_{T_L}$ à valeurs réelles. La description correcte est
la loi asymptotique de $L_n$, étudiée ci-dessous.

\section{Régime stochastique non dégénéré}
\label{sec:asymptotic-stochastic}

\subsection{Mort presque sûre de l'emprunteuse}

\begin{theorem}[Absorption presque sûre]
Supposons $q>0$ et $\sigma>0$. Alors, pour tout
$0\leq\delta<1$, tout $r\in[0,1]$ et tout état initial fini,
\begin{equation}
\Pp(T_B<\infty)=1.
\label{eq:absorption-as}
\end{equation}
Il n'existe donc pas, avec les chocs exacts de M4.2, de régime de survie
éternelle des deux entités avec probabilité positive.
\end{theorem}
\begin{proof}
Couplons l'emprunteuse, jusqu'à sa mort, à la chaîne autarcique
\begin{equation}
U_{n+1}=d\Phi_\gamma(U_nX^B_{n+1}),\qquad U_0=B_0.
\label{eq:upper-chain}
\end{equation}
Par monotonie de $\Phi_\gamma$ et parce que $c\geq0$, on a
$B_n\leq U_n$ tant que l'emprunteuse vit.

Choisissons $p\in(0,1)$. Le choc lognormal vérifie
\begin{equation}
\E[X^p]=\exp\!\left(\frac{p^2-p}{2}\sigma^2\right)<1.
\label{eq:lognormal-moment}
\end{equation}
Comme $(a+b)^p\leq a^p+b^p$,
\begin{equation}
\E[U_{n+1}^p\mid U_n=u]
\leq d^p\E[X^p]u^p
+d^p\E[X^{\gamma p}]u^{\gamma p}.
\label{eq:drift}
\end{equation}
Le coefficient du terme dominant est strictement inférieur à un et
$\gamma p<p$ : la chaîne ne peut donc s'échapper vers $+\infty$. Il faut
aussi exclure une fuite vers zéro. Pour $s>0$, on a
\[
U_{n+1}^{-s}
\leq d^{-s}U_n^{-\gamma s}(X^B_{n+1})^{-\gamma s}.
\]
Tous les moments négatifs du choc lognormal sont finis et
$u^{-\gamma s}/u^{-s}=u^{s(1-\gamma)}\to0$ quand $u\downarrow0$.
Une fonction de Lyapunov combinant $u^p$ et $u^{-s}$ fournit donc une
dérive vers l'intérieur aux deux extrémités. Jointe à la densité
strictement positive du choc, elle implique que la chaîne $U$ revient
presque sûrement une infinité de fois dans un compact
$[\varepsilon,M]\subset(0,\infty)$.

À chacune de ces visites, si l'emprunteuse est encore vivante, alors
$q\leq B_n\leq U_n\leq M$. Le choc suivant la tue dès que
$X^B_{n+1}<z_*/M$, événement de probabilité
$\Pp(X<z_*/M)>0$, uniforme sur ce compact. La propriété de Markov forte
aux temps de retour entraîne que la probabilité d'échapper à toutes ces
tentatives est nulle.
\end{proof}

Le théorème inclut $r=0$. Même sans rente, la dette nominale $q$ subsiste
et la condition de valeur nette impose $B_n\geq q$. Un choc lognormal
arbitrairement défavorable finit presque sûrement par franchir ce seuil.

\subsection{Dynamique de la prêteuse après le défaut}

Après $T_B$, la prêteuse suit la chaîne autarcique
\begin{equation}
L_{n+1}=d\Phi_\gamma(L_nX^L_{n+1}).
\label{eq:lender-autarky}
\end{equation}
Son noyau a pour densité, pour $l,y>0$,
\begin{equation}
A(l,y)=
\frac{\rho_\sigma\!\left(\Psi_\gamma(y/d)/l\right)}
{d\,l\,\Phi_\gamma'\!\left(\Psi_\gamma(y/d)\right)}.
\label{eq:autarky-kernel}
\end{equation}
La même dérive \eqref{eq:drift} assure l'existence d'une unique loi
invariante $\pi_{\gamma,\delta,\sigma}$ et la convergence en loi vers
celle-ci. Sa densité, lorsqu'elle est notée $\pi(y)$, est caractérisée par
l'équation de Fredholm
\begin{equation}
\pi(y)=\int_0^\infty \pi(l)A(l,y)\dd l,
\qquad \int_0^\infty\pi(y)\dd y=1.
\label{eq:invariant-fredholm}
\end{equation}
Elle ne possède pas, en général, de forme élémentaire.

Si $\delta>0$ et le premier moment est fini, la stationnarité donne
l'identité exacte
\begin{equation}
\delta\E_\pi[L]
=d\,e^{-\gamma(1-\gamma)\sigma^2/2}
\E_\pi[L^\gamma].
\label{eq:stationary-moment}
\end{equation}
En effet, $L$ et le nouveau choc sont indépendants sous la loi
stationnaire, $\E X=1$ et
$\E X^\gamma=e^{-\gamma(1-\gamma)\sigma^2/2}$.

Pour $p>0$, le coefficient multiplicatif asymptotique satisfait
\begin{equation}
\E[(dX)^p]
=\exp\!\left(p\log d+\frac{p^2-p}{2}\sigma^2\right).
\label{eq:moment-coefficient}
\end{equation}
Lorsque $\delta>0$, le seuil positif non nul résolvant
$\E[(dX)^p]=1$ est
\begin{equation}
p_*=1-\frac{2\log d}{\sigma^2}>1.
\label{eq:pstar}
\end{equation}
L'inégalité de dérive fournit notamment la finitude des moments d'ordre
$p<p_*$ compatibles avec le terme sous-linéaire. Cette relation explique
comment une dépréciation plus forte allège la queue stationnaire, tandis
qu'une volatilité plus forte l'alourdit.

\begin{corollary}[État asymptotique du système]
Pour $q>0$ et $\sigma>0$, sur l'espace d'états complété par le cimetière
$\partial$,
\begin{equation}
\Law(L_n,B_n)\Longrightarrow
\pi_{\gamma,\delta,\sigma}\otimes\delta_{\partial}.
\label{eq:system-limit}
\end{equation}
La loi limite ne dépend ni de $q$ ni de $r$ ; ces paramètres gouvernent le
transitoire, la loi de $T_B$, les capitaux au défaut et la richesse
accumulée par la prêteuse avant le retour progressif au régime autarcique.
\end{corollary}

Des chocs persistants non dégénérés excluent une convergence presque sûre
vers une constante : même sous la loi invariante, un nouveau multiplicateur
aléatoire intervient à chaque pas. La formulation « tailles constantes
positives » doit donc être réservée à $\sigma=0$.

\section{Régime déterministe \texorpdfstring{$\sigma=0$}{sigma = 0}}

Cette section décrit exactement les cas où les deux tailles peuvent
converger vers des constantes. Supposons d'abord $0<\delta<1$ et posons
\begin{equation}
a=\frac{\delta}{1-\delta}=\frac\delta d.
\label{eq:a}
\end{equation}
Tant que l'emprunteuse vit, sa récurrence est
\begin{equation}
B_{n+1}=G_c(B_n)=d(B_n+B_n^\gamma-c),
\qquad B_{n+1}\geq q.
\label{eq:det-borrower}
\end{equation}
Ses points fixes positifs résolvent
\begin{equation}
H(k)=k^\gamma-ak=c.
\label{eq:H}
\end{equation}
La fonction $H$ est strictement concave et atteint son maximum en
\begin{equation}
k_m=\left(\frac\gamma a\right)^{1/(1-\gamma)},
\qquad
c_{\max}=H(k_m)
=(1-\gamma)\left(\frac\gamma a\right)^{\gamma/(1-\gamma)}.
\label{eq:cmax}
\end{equation}

\begin{theorem}[Classification déterministe de l'emprunteuse]
Soit $B_0>q$.
\begin{enumerate}
\item Si $0\leq c<c_{\max}$, l'équation \eqref{eq:H} possède deux
racines $k_-<k_m<k_+$, avec la convention $k_-=0$ lorsque $c=0$.
La racine $k_-$ est instable et $k_+$ est stable. L'emprunteuse survit
éternellement et converge vers $k_+$ si
\begin{equation}
B_0>k_-,\qquad k_+\geq q.
\label{eq:det-survival-condition}
\end{equation}
Le cas exceptionnel $B_0=k_-\geq q$ donne une trajectoire constante
instable. Dans tous les autres cas, elle franchit $q$ et meurt en temps
fini.
\item Si $c=c_{\max}$, le point $k_m$ est double. La survie a lieu si
$B_0\geq k_m\geq q$ ; la convergence vers $k_m$ est alors non
géométrique. Sinon le décès est fini.
\item Si $c>c_{\max}$, aucun point fixe n'existe et le décès survient en
temps fini.
\end{enumerate}
\end{theorem}
\begin{proof}
On a
\[
G_c(k)-k=d\{H(k)-c\},\qquad
G_c'(k)=d(1+\gamma k^{\gamma-1}).
\]
La géométrie de la fonction strictement concave $H$ donne le nombre de
racines et le signe de $G_c(k)-k$. Au point $k_m$, $G_c'(k_m)=1$ ; la
racine inférieure a donc une dérivée supérieure à un et la racine
supérieure une dérivée inférieure à un. La monotonie de $G_c$ donne les
bassins annoncés. Si la limite attractive est sous $q$, ou si la suite
décroît vers zéro, le seuil $q$ est franchi en temps fini.
\end{proof}

Si l'emprunteuse survit, la prêteuse reçoit indéfiniment $c$ et suit
\[
L_{n+1}=d(L_n+L_n^\gamma+c).
\]
Elle converge vers l'unique solution $\ell_c>K_{\mathrm{aut}}$ de
\begin{equation}
a\ell_c-\ell_c^\gamma=c,
\qquad
K_{\mathrm{aut}}=a^{-1/(1-\gamma)}
=\left(\frac{1-\delta}{\delta}\right)^{1/(1-\gamma)}.
\label{eq:lender-fixed}
\end{equation}
Si l'emprunteuse meurt, la rente cesse et la prêteuse converge vers
$K_{\mathrm{aut}}$.

Pour le cas central $\gamma=1/2$, toutes les frontières sont explicites.
Avec $u=\sqrt{k}$,
\begin{equation}
k_\pm=\left(\frac{1\pm\sqrt{1-4ac}}{2a}\right)^2,
\quad c_{\max}=\frac1{4a},
\quad
\ell_c=\left(\frac{1+\sqrt{1+4ac}}{2a}\right)^2.
\label{eq:sqrt-explicit}
\end{equation}

Si $\delta=0$ et $\sigma=0$, il n'existe pas de force de rappel
autarcique. L'emprunteuse vérifie
$B_{n+1}-B_n=B_n^\gamma-c$. Elle croît sans borne si
$B_0>c^{1/\gamma}$, reste constante à l'égalité, et décroît jusqu'au
défaut si $B_0<c^{1/\gamma}$. La prêteuse croît toujours sans borne. Ce
cas frontière n'admet donc pas les constantes autarciques de
\eqref{eq:lender-fixed}.

\section{Effets comparatifs de
\texorpdfstring{$r$, $\delta$, $\sigma$ et $\gamma$}{r, delta, sigma et gamma}}

\subsection{Taux et dépréciation : ordres trajectoriels}

À $q$ et état initial fixés, un taux plus élevé augmente $c=rq$ et diminue
la récurrence de l'emprunteuse pour toute réalisation commune des chocs.
Par couplage monotone,
\begin{equation}
r_2\geq r_1\quad\Longrightarrow\quad
T_B(r_2)\leq T_B(r_1)\quad\text{presque sûrement}.
\label{eq:r-monotone}
\end{equation}
Une dépréciation plus élevée diminue directement
$d\{\Phi_\gamma(BX)-c\}$ et augmente le seuil
$c+q/d$. De même,
\begin{equation}
\delta_2\geq\delta_1\quad\Longrightarrow\quad
T_B(\delta_2)\leq T_B(\delta_1)
\quad\text{presque sûrement}.
\label{eq:delta-monotone}
\end{equation}
Ces deux classements sont globaux. Le changement simultané de $q$,
$c=rq$ et $B_0$ entre les règles arithmétique et géométrique est, lui,
résolu par la normalisation \eqref{eq:normalized-borrower} :
$T_B^{\mathrm A}\leq T_B^{\mathrm G}$ presque sûrement.

\subsection{Amplitude du choc}

Pour un état courant $b$, écrivons $\alpha=z_*/b$. Le risque instantané
est
\begin{equation}
\kappa_\sigma(b)=
\Phi_{\mathrm N}\left(\frac{\log\alpha}{\sigma}+\frac\sigma2\right).
\label{eq:kappa-sigma}
\end{equation}
La dérivée de l'argument de $\Phi_{\mathrm N}$ vaut
\begin{equation}
\frac12-\frac{\log\alpha}{\sigma^2}.
\label{eq:sigma-derivative}
\end{equation}
Si $\alpha\leq1$, le risque à un pas croît avec $\sigma$. Si
$\alpha>1$, il décroît d'abord puis croît après
$\sigma=\sqrt{2\log\alpha}$. La volatilité n'est donc pas globalement
ordonnée, même au premier pas : le déplacement de la moyenne logarithmique
$-\sigma^2/2$, nécessaire pour conserver $\E X=1$, intervient avec
l'étalement.

\subsection{Concavité}

À échelle fixée,
\[
\frac{\partial}{\partial\gamma}\Phi_\gamma(z)=z^\gamma\log z.
\]
Augmenter $\gamma$ accroît la production si $z>1$, mais la réduit si
$0<z<1$. Il n'existe donc pas d'ordre global en $\gamma$ sans
normalisation d'échelle. Toute comparaison doit passer par
$\Psi_\gamma(h)$ dans \eqref{eq:hazard} ou par le noyau
\eqref{eq:Qdensity}. Ce point est essentiel : un effet attribué à la
concavité peut en réalité provenir du déplacement de l'échelle autarcique
$K_{\mathrm{aut}}$.

\section{Choix du taux selon l'objectif de la prêteuse}

Le modèle mathématique ne sélectionne pas un taux unique lorsque
$r\in[0,1]$ est une décision stratégique. Il faut préciser l'objectif.
Quelques résultats exacts délimitent les possibilités.

\paragraph{Recette personnelle au prochain pas.}
La recette effective est
\begin{equation}
R_1(r)=\E\bigl[\min(rq,\Phi_\gamma(B_0X))\bigr].
\label{eq:immediate-revenue}
\end{equation}
Pour $\sigma>0$, elle possède elle aussi une expression lognormale fermée.
Pour $x>0$ et $p>0$, posons
\begin{equation}
M_p(x)=\E[X^p\1_{\{X<x\}}]
=e^{(p^2-p)\sigma^2/2}
\Phi_{\mathrm N}\!\left(
\frac{\log x+(1/2-p)\sigma^2}{\sigma}\right).
\label{eq:truncated-moment}
\end{equation}
Si $c=rq>0$ et $x_c=\Psi_\gamma(c)/B_0$, alors
\begin{equation}
R_1(r)=B_0M_1(x_c)+B_0^\gamma M_\gamma(x_c)
+c\{1-\Pp(X<x_c)\}.
\label{eq:closed-revenue}
\end{equation}
Pour $c=0$, $R_1=0$.
Elle est croissante en $r$ point par point. Une prêteuse maximisant
uniquement sa recette immédiate choisit donc $r=1$. La même conclusion
vaut pour son capital espéré immédiatement après dépréciation.

\paragraph{Survie de l'emprunteuse.}
La probabilité de survie à tout horizon décroît avec $r$ par
\eqref{eq:r-monotone}. Une prêteuse qui souhaite uniquement préserver son
partenaire choisit $r=0$. Cela ne garantit toutefois pas une survie
éternelle lorsque $\sigma>0$, à cause de la contrainte de valeur nette
$B\geq q$.

\paragraph{Production commune instantanée.}
Le taux n'agit pas sur la production du pas d'émission. Le principal
$q_{\mathrm A}$ est optimal par stricte concavité. Aux pas ultérieurs, le
service transfère du capital de l'entité pauvre vers l'entité riche et
tend à réduire la production concave commune tant que les deux survivent.
La faillite ajoute une discontinuité par destruction du capital résiduel ;
aucun taux constant n'est universellement optimal.

\paragraph{Objectif intertemporel de la prêteuse.}
Pour un facteur d'actualisation $\beta\in(0,1)$ et une utilité $U$, une
formulation exacte est
\begin{equation}
V_r(l,b)=\E_{l,b}\left[\sum_{n\geq0}\beta^nU(L_n)\right],
\qquad r\in[0,1],
\label{eq:bellman-objective}
\end{equation}
où les transitions avant et au défaut sont données par
$\mathcal S$ et $\mathcal D$, puis par $A$. Le choix optimal est
$\arg\max_{r\in[0,1]}V_r(L_0,B_0)$. La rente croissante et la durée de
vie décroissante créent un arbitrage réel ; sans spécifier $U$ et $\beta$,
il est mathématiquement impossible de sélectionner un taux intérieur.

\section{Synthèse des issues asymptotiques}

\begin{center}
\begin{tabularx}{\textwidth}{@{}lXX@{}}
\toprule
Régime & Emprunteuse & Prêteuse \\
\midrule
$\sigma>0$, $q>0$
& $T_B<\infty$ p.s.; loi exacte par
$\delta_{B_0}Q^nD$; puis état $\partial$
& $T_L=\infty$; loi au défaut par
$(I-\mathcal S)^{-1}\mathcal D$; puis convergence en loi vers $\pi$ \\
\addlinespace
$\sigma=0$, $0<\delta<1$, survie de $B$
& convergence vers $k_+$ sous les conditions
\eqref{eq:det-survival-condition}
& convergence vers $\ell_c$ \\
\addlinespace
$\sigma=0$, $0<\delta<1$, mort de $B$
& mort déterministe en temps fini
& convergence vers $K_{\mathrm{aut}}$ \\
\addlinespace
$\sigma=\delta=0$
& croissance, équilibre instable ou mort selon $B_0^\gamma-c$
& croissance sans borne \\
\bottomrule
\end{tabularx}
\end{center}

Les deux issues envisagées « mort de l'une puis de l'autre » et
« prêteuse morte, emprunteuse survivante » sont exclues par le bilan exact
de cette expérience. Sous choc non dégénéré, l'issue presque sûre est
unique au niveau qualitatif : mort finie de l'emprunteuse, survie éternelle
de la prêteuse. Les paramètres et la règle de principal restent néanmoins
déterminants pour les distributions de $T_B$, $B_T^{\cim}$ et $L_{T_B}$.

\section{Calcul effectif sans simulation}

Pour calculer numériquement les lois théoriques sans simuler de
trajectoires, une discrétisation déterministe du noyau suffit :
\begin{enumerate}
\item calculer $q_s,L_0^s,B_0^s$ pour
$s\in\{\mathrm A,\mathrm G\}$ ;
\item résoudre scalairement
$z+z^\gamma=q_s(r+1/d)$ pour obtenir $z_*^s$ ;
\item évaluer \eqref{eq:hazard} pour la probabilité au premier pas ;
\item discrétiser $b,y\in[q_s,\infty)$ sur une grille transformée en
$\log b$, puis construire la matrice positive correspondant à
\eqref{eq:Qdensity} ;
\item itérer $\mu_{n+1}=\mu_nQ$ et utiliser
\eqref{eq:Tlaw}, \eqref{eq:joint-T-B} ;
\item pour la loi jointe avec la prêteuse, appliquer le noyau produit
\eqref{eq:joint-survival-kernel} et le noyau terminal
\eqref{eq:joint-death-kernel}.
\end{enumerate}
Cette procédure est une quadrature des formules analytiques, non une
simulation de Monte-Carlo. Un code empirique minimal séparé pourra ensuite
contrôler ces distributions, en particulier le risque à un pas et les
relations d'ordre en $r$ et $\delta$.

\section{Conclusion}

L'expérience à deux entités se réduit exactement à une chaîne de Markov
tuée unidimensionnelle pour l'emprunteuse, couplée à la prêteuse seulement
par le paiement terminal. La totalité des lois d'arrêt est encodée par le
noyau explicite $Q_{q,r}$ et sa résolvante. Le taux agit à travers la rente
$c=rq$ ; la dépréciation agit à la fois sur la dynamique et sur le seuil
de solvabilité $q/(1-\delta)$ ; la volatilité intervient par une fonction
normale fermée au risque conditionnel ; la concavité intervient par
l'inverse de $z+z^\gamma$.

La moyenne arithmétique est l'optimum de production commune au moment du
transfert, mais elle impose plus de principal et plus de levier que la
cible géométrique. Pour la durée de survie, ce coût nominal domine : sous
un couplage commun, $T_B^{\mathrm A}\leq T_B^{\mathrm G}$ presque
sûrement. Enfin, le réseau minimal révèle un fait institutionnel fort :
une créance perdue ne peut
tuer une prêteuse sans dette, car la perte du principal nominal ne rend pas
son capital réel négatif. Les cascades à plusieurs morts observées dans le
modèle complet sont donc nécessairement un effet de réseau, absent de la
dyade isolée.

\end{document}
